% GENERATED FILE. Edit canonical lessons, then run npm run generate:books.
% canonical-source-hash: fnv1a64:06144c8b2a1d0e27
% canonical-lessons: FR-C381-autoriser, FR-C381-autorisation, FR-C381-tolerer, FR-C381-tolerance, FR-C381-exiger, FR-R381-revision-374-381

\chapter{Permis ou non}
\label{ch:fr-permis-ou-non}
\hlchaptermodality{381}

\begin{chapteropening}
\textbf{By the end of this chapter:} \emph{I can say what is authorised or tolerated, talk about permission and tolerance, and say what someone demands.}

Complete the last lesson of chapter 381: I can say what is authorised or tolerated, talk about permission and tolerance, and say what someone demands.
\end{chapteropening}

\section[autoriser]{autoriser — to authorise, to allow}

\label{lesson:FR-C381-autoriser}



\begin{quote}
Before the new one: say the French for optional, then the French for the obligation.
\end{quote}

\subsection*{You'll want to know: autoriser}
\textbf{autoriser} — \textquotedblleft{}to authorise, to allow\textquotedblright{}.

\begin{quote}
Les chiens ne sont pas autorisés dans le magasin.
\end{quote}

\begin{quote}
\emph{Dogs are not allowed in the shop.}
\end{quote}

\emph{Autoriser quelqu'un à} + infinitive.

\begin{cousinweb}
\emph{autoriser} comes from Old French \emph{auctorisier} (EtymDB), from Medieval Latin \emph{auctōrizāre}, built on \emph{auctor} \textquotedblleft{}author, guarantor\textquotedblright{}, from \emph{augēre} \textquotedblleft{}to increase\textquotedblright{} (Lewis \& Short).
\end{cousinweb}

\begin{grammarlens}[title={What you've built}]
One more word for telling, explaining and describing at B1.
\end{grammarlens}

\subsection*{Your turn}
\begin{itemize}
\raggedright
  \item \emph{Say it:} \emph{autoriser}
  \item \emph{Say it:} \emph{autoriser}, once more
  \item \emph{Say it:} a sentence of your own with \emph{autoriser}
  \item \emph{Recall:} say the French for optional, then the French for the obligation, then say \emph{autoriser} again
\end{itemize}

\begin{tcolorbox}[breakable,colback=teal!4,colframe=teal!35!black,title={Before you move on}]
What does autoriser mean? (\textquotedblleft{}To authorise, to allow\textquotedblright{}.) Say it once more.
\end{tcolorbox}
\section[l'autorisation]{l'autorisation — the authorisation, the permission}

\label{lesson:FR-C381-autorisation}



\begin{quote}
Before the new one: say the French for the obligation, then the French for to authorise.
\end{quote}

\subsection*{You'll want to know: l'autorisation}
\textbf{l'autorisation} — \textquotedblleft{}the authorisation, the permission\textquotedblright{}.

\begin{quote}
Il faut une autorisation pour filmer ici.
\end{quote}

\begin{quote}
\emph{You need permission to film here.}
\end{quote}

Feminine. Close to \emph{la permission}, but more official.

\begin{cousinweb}
\emph{autorisation} is built from \emph{autoriser} and \emph{-ation}. The verb rests on Latin \emph{auctor} \textquotedblleft{}author, guarantor\textquotedblright{}, from \emph{augēre} \textquotedblleft{}to increase\textquotedblright{} (Lewis \& Short).
\end{cousinweb}

\begin{grammarlens}[title={What you've built}]
One more word for telling, explaining and describing at B1.
\end{grammarlens}

\subsection*{Your turn}
\begin{itemize}
\raggedright
  \item \emph{Say it:} \emph{l'autorisation}
  \item \emph{Say it:} \emph{l'autorisation}, once more
  \item \emph{Say it:} a sentence of your own with \emph{l'autorisation}
  \item \emph{Recall:} say the French for the obligation, then the French for to authorise, then say \emph{l'autorisation} again
\end{itemize}

\begin{tcolorbox}[breakable,colback=teal!4,colframe=teal!35!black,title={Before you move on}]
What does l'autorisation mean? (\textquotedblleft{}The authorisation, the permission\textquotedblright{}.) Say it once more.
\end{tcolorbox}
\section[tolérer]{tolérer — to tolerate, to put up with}

\label{lesson:FR-C381-tolerer}



\begin{quote}
Before the new one: say the French for to authorise, then the French for the authorisation.
\end{quote}

\subsection*{You'll want to know: tolérer}
\textbf{tolérer} — \textquotedblleft{}to tolerate, to put up with\textquotedblright{}.

\begin{quote}
Le professeur ne tolère pas les retards.
\end{quote}

\begin{quote}
\emph{The teacher does not tolerate lateness.}
\end{quote}

\emph{Je tolère}, with \emph{è} before a silent ending.

\begin{cousinweb}
\emph{tolérer} is a learned borrowing from Latin \emph{tolerāre} \textquotedblleft{}to bear, to endure\textquotedblright{} (EtymDB; Lewis \& Short).
\end{cousinweb}

\begin{grammarlens}[title={What you've built}]
One more word for telling, explaining and describing at B1.
\end{grammarlens}

\subsection*{Your turn}
\begin{itemize}
\raggedright
  \item \emph{Say it:} \emph{tolérer}
  \item \emph{Say it:} \emph{tolérer}, once more
  \item \emph{Say it:} a sentence of your own with \emph{tolérer}
  \item \emph{Recall:} say the French for to authorise, then the French for the authorisation, then say \emph{tolérer} again
\end{itemize}

\begin{tcolorbox}[breakable,colback=teal!4,colframe=teal!35!black,title={Before you move on}]
What does tolérer mean? (\textquotedblleft{}To tolerate, to put up with\textquotedblright{}.) Say it once more.
\end{tcolorbox}
\section[la tolérance]{la tolérance — the tolerance}

\label{lesson:FR-C381-tolerance}



\begin{quote}
Before the new one: say the French for the authorisation, then the French for to tolerate.
\end{quote}

\subsection*{You'll want to know: la tolérance}
\textbf{la tolérance} — \textquotedblleft{}the tolerance\textquotedblright{}.

\begin{quote}
Il faut un peu de tolérance pour vivre ensemble.
\end{quote}

\begin{quote}
\emph{Living together takes a bit of tolerance.}
\end{quote}

\emph{Tolérance zéro} — zero tolerance.

\begin{cousinweb}
\emph{tolérance} is a learned borrowing from Latin \emph{tolerantia} \textquotedblleft{}endurance, bearing\textquotedblright{} (Lewis \& Short), from \emph{tolerāre} \textquotedblleft{}to bear\textquotedblright{}.
\end{cousinweb}

\begin{grammarlens}[title={What you've built}]
One more word for telling, explaining and describing at B1.
\end{grammarlens}

\subsection*{Your turn}
\begin{itemize}
\raggedright
  \item \emph{Say it:} \emph{la tolérance}
  \item \emph{Say it:} \emph{la tolérance}, once more
  \item \emph{Say it:} a sentence of your own with \emph{la tolérance}
  \item \emph{Recall:} say the French for the authorisation, then the French for to tolerate, then say \emph{la tolérance} again
\end{itemize}

\begin{tcolorbox}[breakable,colback=teal!4,colframe=teal!35!black,title={Before you move on}]
What does la tolérance mean? (\textquotedblleft{}The tolerance\textquotedblright{}.) Say it once more.
\end{tcolorbox}
\section[exiger]{exiger — to demand, to require}

\label{lesson:FR-C381-exiger}



\begin{quote}
Before the new one: say the French for to tolerate, then the French for the tolerance.
\end{quote}

\subsection*{You'll want to know: exiger}
\textbf{exiger} — \textquotedblleft{}to demand, to require\textquotedblright{}.

\begin{quote}
Le patron exige beaucoup de ses salariés.
\end{quote}

\begin{quote}
\emph{The boss demands a lot of his employees.}
\end{quote}

\emph{Exiger que} takes the subjunctive. \emph{Nous exigeons}.

\begin{cousinweb}
\emph{exiger} is a learned borrowing from Latin \emph{exigere} \textquotedblleft{}to drive out, to demand\textquotedblright{} (Lewis \& Short), built from \emph{ex} \textquotedblleft{}out\textquotedblright{} and \emph{agere} \textquotedblleft{}to drive\textquotedblright{}.
\end{cousinweb}

\begin{grammarlens}[title={What you've built}]
One more word for telling, explaining and describing at B1.
\end{grammarlens}

\subsection*{Your turn}
\begin{itemize}
\raggedright
  \item \emph{Say it:} \emph{exiger}
  \item \emph{Say it:} \emph{exiger}, once more
  \item \emph{Say it:} a sentence of your own with \emph{exiger}
  \item \emph{Recall:} say the French for to tolerate, then the French for the tolerance, then say \emph{exiger} again
\end{itemize}

\begin{tcolorbox}[breakable,colback=teal!4,colframe=teal!35!black,title={Before you move on}]
What does exiger mean? (\textquotedblleft{}To demand, to require\textquotedblright{}.) Say it once more.
\end{tcolorbox}
\section[(dialogue)]{Retour sur les chapitres 374–381}

\label{lesson:FR-R381-revision-374-381}



\begin{quote}
Say the French for the aim and the intention.
\end{quote}

\subsection*{Your turn}
Say these aloud:

\begin{itemize}
\raggedright
  \item the aim, the intention, the objective, the resolution, to aim
  \item the secondary school, the secondary-school pupil, to attend, the school report, the head teacher
  \item the faculty, to enrol, the registration, the semester, the lecture hall
  \item to revise, the revision, to memorise, by heart, the textbook
  \item to fail, to repeat a year, the competitive exam, the average, the resit
  \item the diploma, the test, the distinction, the doctorate, the grant
  \item the norm, the regulations, compulsory, optional, the obligation
  \item to authorise, the authorisation, to tolerate, the tolerance, to demand
\end{itemize}

\begin{tcolorbox}[breakable,colback=teal!4,colframe=teal!35!black,title={Before you move on}]
The aim? (\textbf{le but}.) The intention? (\textbf{l'intention}.) The objective? (\textbf{l'objectif}.) The resolution? (\textbf{la résolution}.) To aim? (\textbf{viser}.) The secondary school? (\textbf{le lycée}.) The secondary-school pupil? (\textbf{le lycéen}.) To attend? (\textbf{fréquenter}.) The school report? (\textbf{le bulletin}.) The head teacher? (\textbf{le proviseur}.) The faculty? (\textbf{la faculté}.) To enrol? (\textbf{inscrire}.) The registration? (\textbf{l'inscription}.) The semester? (\textbf{le semestre}.) The lecture hall? (\textbf{l'amphithéâtre}.) To revise? (\textbf{réviser}.) The revision? (\textbf{la révision}.) To memorise? (\textbf{mémoriser}.) By heart? (\textbf{par cœur}.) The textbook? (\textbf{le manuel}.) To fail? (\textbf{échouer}.) To repeat a year? (\textbf{redoubler}.) The competitive exam? (\textbf{le concours}.) The average? (\textbf{la moyenne}.) The resit? (\textbf{le rattrapage}.) The diploma? (\textbf{le diplôme}.) The test? (\textbf{l'épreuve}.) The distinction? (\textbf{la mention}.) The doctorate? (\textbf{le doctorat}.) The grant? (\textbf{la bourse}.) The norm? (\textbf{la norme}.) The regulations? (\textbf{le règlement}.) Compulsory? (\textbf{obligatoire}.) Optional? (\textbf{facultatif}.) The obligation? (\textbf{l'obligation}.) To authorise? (\textbf{autoriser}.) The authorisation? (\textbf{l'autorisation}.) To tolerate? (\textbf{tolérer}.) The tolerance? (\textbf{la tolérance}.) To demand? (\textbf{exiger}.)
\end{tcolorbox}
